%% ============================================================
%%  Quantum Computing and Quantum Machine Learning
%%  Week 15: QNN, VQE, QAOA, QSVM, HHL
%%  Author: Morten Hjorth-Jensen
%%  Department of Physics, University of Oslo
%% ============================================================
\documentclass[aspectratio=169,10pt]{beamer}

%% ── Theme & colours ─────────────────────────────────────────────
\usetheme{Madrid}
\usecolortheme{seahorse}
\usefonttheme[onlymath]{serif}

\setbeamercolor{frametitle}{bg=blue!15!white,fg=blue!60!black}
\setbeamercolor{title}{fg=blue!70!black}
\setbeamercolor{block title}{bg=blue!20!white,fg=blue!60!black}
\setbeamercolor{block body}{bg=blue!5!white}
\setbeamercolor{alerted text}{fg=red!60!black}
\setbeamertemplate{navigation symbols}{}
\setbeamertemplate{footline}{%
  \leavevmode\hbox{%
    \begin{beamercolorbox}[wd=.25\paperwidth,ht=2.25ex,dp=1ex,center]{author in head/foot}%
      \usebeamerfont{author in head/foot}\insertshortauthor
    \end{beamercolorbox}%
    \begin{beamercolorbox}[wd=.50\paperwidth,ht=2.25ex,dp=1ex,center]{title in head/foot}%
      \usebeamerfont{title in head/foot}\insertshorttitle
    \end{beamercolorbox}%
    \begin{beamercolorbox}[wd=.25\paperwidth,ht=2.25ex,dp=1ex,right]{date in head/foot}%
      \usebeamerfont{date in head/foot}\insertshortdate{}\hspace*{1em}%
      \insertframenumber{}/\inserttotalframenumber\hspace*{1ex}
    \end{beamercolorbox}}%
  \vskip0pt}

%% ── Packages ────────────────────────────────────────────────────
\usepackage{amsmath,amssymb,bm,physics}
\usepackage{braket}
\usepackage{graphicx}
\usepackage{booktabs}
\usepackage{xcolor}
\usepackage{tikz}
\usepackage{quantikz}
\usepackage{listings}
\usepackage{hyperref}

%% ── Python listing style ─────────────────────────────────────────
\lstset{
  language=Python,
  basicstyle=\ttfamily\scriptsize,
  keywordstyle=\color{blue!70!black}\bfseries,
  commentstyle=\color{gray},
  stringstyle=\color{red!60!black},
  breaklines=true,
  frame=single,
  backgroundcolor=\color{blue!3!white},
  xleftmargin=4pt, xrightmargin=4pt
}

%% ── Metadata ────────────────────────────────────────────────────
\title[QC and QML -- Week 14]{Quantum Computing and Quantum Machine Learning}
\subtitle{QNN, VQE, QAOA, QSVM, HHL}
\author[M.~Hjorth-Jensen]{Morten Hjorth-Jensen}
\institute[UiO]{Department of Physics \& CCSE\\University of Oslo}
\date{April 29, 2026}

%% ============================================================
\begin{document}
%% ============================================================

\begin{frame}
  \titlepage
\end{frame}

\begin{frame}{Outline}
  \tableofcontents[hidesubsections]
\end{frame}

%% ============================================================
\section{Introduction to Quantum Machine Learning}
%% ============================================================

\begin{frame}{Week 14 — Plan}
  \begin{enumerate}\itemsep4pt
    \item Discussion of the \textbf{QAOA algorithm} with repetition from last week
    \item Parametrized quantum circuits (PQC) and Variational Quantum Circuits (VQCs)
    \item \textbf{Quantum Neural Networks (QNNs):} mathematical structure, QFI, TDVP,
          barren plateaus, quantikz circuit examples
    \item \textbf{Variational Quantum Eigensolver (VQE):} variational ansatz, UCC,
          connections to many-body theory
    \item \textbf{QAOA:} adiabatic motivation, MaxCut, parameter landscape
    \item \textbf{QSVM:} quantum kernels, fidelity, RKHS
    \item \textbf{HHL algorithm:} solving linear systems, QPE, Green's functions
    \item Grand unification: Geometry + Dynamics + Operators
  \end{enumerate}
\end{frame}

\begin{frame}{What is Quantum Machine Learning?}
  \begin{columns}[T]
    \column{0.55\textwidth}
      Quantum Machine Learning (QML) integrates quantum computing with
      machine learning to exploit quantum advantages:
      \begin{itemize}
        \item \textbf{Quantum parallelism:} process exponentially many states
        \item \textbf{Quantum entanglement:} correlated, non-classical feature spaces
        \item \textbf{Quantum interference:} enhance solutions, suppress errors
      \end{itemize}
      \vspace{0.3cm}
      \begin{alertblock}{Core question}
        When and how does quantum hardware give a genuine advantage
        over the best classical algorithm?
      \end{alertblock}
    \column{0.42\textwidth}
      \begin{block}{NISQ challenges}
        \begin{itemize}
          \item Noisy Intermediate-Scale Quantum devices
          \item Decoherence limits circuit depth
          \item Efficient classical data encoding is non-trivial
          \item Scalability to large datasets is open
        \end{itemize}
      \end{block}
  \end{columns}
\end{frame}

%% ============================================================
\section{Quantum Neural Networks}
%% ============================================================

\begin{frame}{QNN as a Variational Quantum State}
  The most compact mathematical description of a QNN:
  \[
    \ket{\psi(x,\theta)} = U(\theta)\,U(x)\,\ket{0}^{\otimes n}
  \]
  \[
    f(x,\theta) = \bra{\psi(x,\theta)}\,O\,\ket{\psi(x,\theta)}
  \]
  \begin{columns}[T]
    \column{0.48\textwidth}
      \begin{block}{Three components}
        \begin{itemize}
          \item $U(x)$: \textbf{feature map} — encodes classical data
          \item $U(\theta)$: \textbf{variational ansatz} — trainable unitary
          \item $O$: \textbf{observable} — defines the scalar output
        \end{itemize}
      \end{block}
    \column{0.48\textwidth}
      \begin{alertblock}{Key point}
        Training a QNN means finding $\theta^\star$ that minimises a
        classical loss $\mathcal{L}(\theta)$ computed from repeated
        quantum measurements.
      \end{alertblock}
  \end{columns}
\end{frame}

%% ── Circuit: generic 2-qubit QNN ────────────────────────────────
\begin{frame}{Generic 2-Qubit QNN Circuit}
  \[
    \ket{\psi(x,\theta)} = W(\theta)\,U(x)\,\ket{00}
  \]

  \vspace{0.3cm}
  \begin{center}
  \begin{quantikz}[row sep=0.5cm, column sep=0.5cm]
    \lstick{$\ket{0}$}
      & \gate{R_x(x_1)} & \gate{R_y(\theta_1)} & \ctrl{1}
      & \gate{R_z(\theta_3)} & \meter{} \\
    \lstick{$\ket{0}$}
      & \gate{R_x(x_2)} & \gate{R_y(\theta_2)} & \targ{}
      & \gate{R_z(\theta_4)} & \meter{}
  \end{quantikz}
  \end{center}

  \vspace{0.3cm}
  \begin{itemize}
    \item \textbf{Layer 1} ($U(x)$): $R_x(x_i)$ encodes data as rotation angles
    \item \textbf{Layer 2} ($W_1(\theta)$): $R_y(\theta_i)$ variational rotations
          $\to$ CNOT entangler
    \item \textbf{Layer 3} ($W_2(\theta)$): $R_z(\theta_i)$ refines the state
    \item Output: $f = \langle Z_0 \rangle$ (expectation of $Z\otimes I$)
  \end{itemize}
\end{frame}

%% ── Feature map as Hamiltonian evolution ────────────────────────
\begin{frame}{Feature Map as Hamiltonian Evolution}
  \begin{columns}[T]
    \column{0.52\textwidth}
      A common, physically motivated feature map:
      \[
        U(x) = e^{-iH(x)},\quad
        H(x) = \sum_i x_i Z_i + \sum_{i<j} x_i x_j Z_i Z_j
      \]
      \begin{itemize}
        \item Data $\equiv$ coupling constants of a spin Hamiltonian
        \item Circuit $\equiv$ quantum time evolution
        \item Pairwise terms create entanglement between features
      \end{itemize}
    \column{0.44\textwidth}
      \begin{center}
      \begin{quantikz}[row sep=0.5cm, column sep=0.45cm]
        \lstick{$\ket{0}$}
          & \gate{H} & \gate{R_z(x_1)} & \ctrl{1} & \qw \\
        \lstick{$\ket{0}$}
          & \gate{H} & \gate{R_z(x_2)} & \gate{R_{zz}(x_1 x_2)} & \qw
      \end{quantikz}
      \end{center}
      \vspace{0.2cm}
      {\small $R_{zz}(\phi)=e^{-i\phi Z\otimes Z/2}$:
       two-qubit entangling gate encoding pairwise product $x_1 x_2$.}
  \end{columns}
\end{frame}

%% ── Pauli ansatz ─────────────────────────────────────────────────
\begin{frame}{Pauli Expansion of the Variational Ansatz}
  \begin{columns}[T]
    \column{0.52\textwidth}
      The variational unitary is generated by a Pauli Hamiltonian:
      \[
        U(\theta) = e^{-iH(\theta)},\quad
        H(\theta) = \sum_\alpha \theta_\alpha P_\alpha
      \]
      where $P_\alpha \in \{I,X,Y,Z\}^{\otimes n}$ are Pauli strings.

      \medskip
      The BCH expansion shows what the circuit computes:
      \[
        U^\dagger O\, U = O + i[H,O]
          + \frac{i^2}{2}[H,[H,O]] + \cdots
      \]
      Each commutator level adds higher-body correlations — analogous
      to coupled-cluster expansions.
    \column{0.44\textwidth}
      \begin{center}
      \begin{quantikz}[row sep=0.5cm, column sep=0.4cm]
        \lstick{$\ket{0}$}
          & \gate{R_y(\theta_1)} & \ctrl{1}
          & \gate{R_z(\phi_1)}   & \qw \\
        \lstick{$\ket{0}$}
          & \gate{R_y(\theta_2)} & \targ{}
          & \gate{R_z(\phi_2)}   & \qw \\
        \lstick{$\ket{0}$}
          & \gate{R_y(\theta_3)} & \ctrl{1}
          & \gate{R_z(\phi_3)}   & \qw \\
        \lstick{$\ket{0}$}
          & \gate{R_y(\theta_4)} & \targ{}
          & \gate{R_z(\phi_4)}   & \qw
      \end{quantikz}
      \end{center}
      {\small One layer of a hardware-efficient ansatz (4 qubits):
       $R_y$ rotations + CNOT ladder + $R_z$ rotations.}
  \end{columns}
\end{frame}

%% ── Deep multi-layer QNN circuit ────────────────────────────────
\begin{frame}{Multi-Layer QNN: Full Circuit Diagram}
  \begin{center}
  \begin{quantikz}[row sep=0.45cm, column sep=0.35cm]
    \lstick{$\ket{0}$}
      & \gate[4]{U(x)}
      & \gate{R_y(\theta_1^{(1)})} & \ctrl{1} & \qw
      & \gate{R_y(\theta_1^{(2)})} & \ctrl{1} & \qw
      & \meter{} \\
    \lstick{$\ket{0}$}
      &
      & \gate{R_y(\theta_2^{(1)})} & \targ{}  & \ctrl{1}
      & \gate{R_y(\theta_2^{(2)})} & \targ{}  & \ctrl{1}
      & \qw \\
    \lstick{$\ket{0}$}
      &
      & \gate{R_y(\theta_3^{(1)})} & \qw      & \targ{}
      & \gate{R_y(\theta_3^{(2)})} & \qw      & \targ{}
      & \qw \\
    \lstick{$\ket{0}$}
      & \qw
      & \gate{R_y(\theta_4^{(1)})} & \qw & \qw
      & \gate{R_y(\theta_4^{(2)})} & \qw & \qw
      & \qw
  \end{quantikz}
  \end{center}
  \vspace{0.2cm}
  {\small Feature map $U(x)$ encodes data ($n=4$ qubits); $L=2$ repeated variational layers
   $W^{(l)}(\theta^{(l)})$ build up correlations via $R_y$ rotations and CNOT
   ladders; qubit 0 is measured to give the scalar output $f = \langle Z_0\rangle$.}
\end{frame}

%% ── Gradient: parameter shift ───────────────────────────────────
\begin{frame}{Parameter-Shift Rule: Exact Quantum Gradients}
  \begin{columns}[T]
    \column{0.52\textwidth}
      For a gate $e^{-i\theta P/2}$ generated by a Pauli $P$
      (eigenvalues $\pm\tfrac{1}{2}$), the gradient is \emph{exact}:
      \[
        \frac{\partial f}{\partial\theta}
        = \frac{1}{2}\bigl[f(\theta+\tfrac{\pi}{2})
                          - f(\theta-\tfrac{\pi}{2})\bigr]
      \]
      \begin{itemize}
        \item Requires only 2 circuit evaluations per parameter
        \item Exact (not a finite-difference approximation)
        \item Hardware-friendly: no ancillas needed
      \end{itemize}
      Gradient derivation from commutator:
      \[
        \partial_\theta f
        = \bra{\psi}\,i[H_\theta,\,O]\,\ket{\psi}
      \]
      \textbf{Physical meaning:} gradients are quantum response
      functions.
    \column{0.44\textwidth}
      \begin{center}
      \begin{quantikz}[row sep=0.5cm, column sep=0.4cm]
        \lstick{$\ket{0}$}
          & \gate{U(x)} & \gate{e^{-i(\theta\pm\pi/2)P/2}}
          & \gate{V} & \meter{}
      \end{quantikz}
      \end{center}
      \vspace{0.3cm}
      \begin{alertblock}{Two-point formula}
        Evaluate the circuit at $\theta+\frac{\pi}{2}$ and
        $\theta-\frac{\pi}{2}$; their half-difference gives the exact
        partial derivative.
      \end{alertblock}
  \end{columns}
\end{frame}

%% ── QFI ─────────────────────────────────────────────────────────
\begin{frame}{Quantum Fisher Information Matrix}
  The \textbf{Quantum Fisher Information} (QFI) defines the Riemannian
  metric on the manifold of variational states:
  \[
    F_{ij} = 4\,\mathrm{Re}\!\left(
      \langle\partial_i\psi|\partial_j\psi\rangle
      - \langle\partial_i\psi|\psi\rangle
        \langle\psi|\partial_j\psi\rangle
    \right)
  \]

  \begin{columns}[T]
    \column{0.50\textwidth}
      \begin{block}{Information-geometric metric}
        \[
          ds^2 = \sum_{ij} F_{ij}\,d\theta_i\,d\theta_j
        \]
        Measures the \emph{distinguishability} of nearby quantum
        states — the Fubini-Study metric on projective Hilbert space.
      \end{block}
      \vspace{0.2cm}
      \begin{block}{Practical formula (overlap method)}
        \[
          F_{ij} = 4\,\mathrm{Re}\left[
            \bra{\partial_i\psi}\!\left(I - \ket{\psi}\bra{\psi}\right)
            \!\ket{\partial_j\psi}
          \right]
        \]
      \end{block}
    \column{0.46\textwidth}
      \begin{alertblock}{Natural gradient descent}
        The \emph{natural gradient} corrects the Euclidean gradient
        for the curved geometry of Hilbert space:
        \[
          \dot{\theta} = F^{-1}\nabla_\theta\mathcal{L}
        \]
        This is the basis of \textbf{Stochastic Reconfiguration (SR)}
        in VMC and of imaginary-time TDVP.
      \end{alertblock}
      \vspace{0.2cm}
      Diagonal elements: $F_{ii}$ = variance of the Pauli generator
      $P_i$ in state $\ket{\psi}$.
  \end{columns}
\end{frame}

%% ── QFI circuit ─────────────────────────────────────────────────
\begin{frame}{Computing the QFI: Hadamard-Test Circuit}
  The overlap $\langle\partial_i\psi|\partial_j\psi\rangle$ entering
  $F_{ij}$ is measured by a Hadamard-test circuit:

  \vspace{0.2cm}
  \begin{center}
  \begin{quantikz}[row sep=0.5cm, column sep=0.5cm]
    \lstick{$\ket{0}$\,(ancilla)}
      & \gate{H} & \ctrl{1} & \gate{H} & \meter{} \\
    \lstick{$\ket{0}^{\otimes n}$}
      & \gate[1]{U(\theta)} & \gate{P_i}
      & \gate[1]{U^\dagger(\theta)} & \qw
  \end{quantikz}
  \end{center}

  \vspace{0.3cm}
  Outcome probability:
  \[
    P(0) = \frac{1}{2}
      + \frac{1}{2}\,\mathrm{Re}\langle\psi|P_i|\psi\rangle
    \implies
    F_{ii} = 4\,\bigl(
      \langle\psi|P_i^2|\psi\rangle - \langle\psi|P_i|\psi\rangle^2
    \bigr)
    = 4\,\mathrm{Var}(P_i)
  \]

  \begin{itemize}
    \item For off-diagonal elements $F_{ij}$ use two Hadamard tests
          with $P_i \pm P_j$
    \item The QFI upper-bounds the classical Fisher information by the
          quantum Cramér-Rao inequality: $\mathrm{Var}(\hat\theta)\ge 1/F$
  \end{itemize}
\end{frame}

%% ── TDVP ─────────────────────────────────────────────────────────
\begin{frame}{Time-Dependent Variational Principle (TDVP)}
  \begin{columns}[T]
    \column{0.52\textwidth}
      The TDVP projects the Schrödinger equation onto the tangent
      space of the variational manifold:
      \[
        \delta\|(i\partial_t - H)|\psi\rangle\| = 0
      \]
      Projecting $\langle\partial_i\psi|(i\partial_t-H)|\psi\rangle=0$:
      \[
        \sum_j F_{ij}\,\dot\theta_j = C_i,\quad
        C_i = \mathrm{Im}\langle\partial_i\psi|H|\psi\rangle
      \]
    \column{0.44\textwidth}
      \begin{block}{Interpretation}
        QNN training $\equiv$ classical dynamical system on a curved
        manifold.  Parameters $\theta_j$ evolve like
        generalised coordinates; $F_{ij}$ is the mass matrix.
      \end{block}
      \vspace{0.2cm}
      \begin{alertblock}{Imaginary-time TDVP}
        Replace $t\to-i\tau$ (imaginary time):
        \[
          \sum_j F_{ij}\,\dot\theta_j = E_i
        \]
        This is \textbf{Stochastic Reconfiguration} —
        optimal ground-state search on the variational manifold.
      \end{alertblock}
  \end{columns}
\end{frame}

%% ── Barren plateaus ─────────────────────────────────────────────
\begin{frame}{Barren Plateaus}
  \begin{columns}[T]
    \column{0.52\textwidth}
      The gradient variance scales exponentially:
      \[
        \mathrm{Var}(\nabla_\theta \mathcal{L}) \sim \frac{1}{2^n}
      \]
      \textbf{Why?} Random, deep circuits form approximate
      \textbf{unitary 2-designs} — their output distribution
      concentrates near the Haar measure.  Local observables
      concentrate around their mean, making gradients vanish.

      \medskip
      \textbf{Strategies to mitigate:}
      \begin{itemize}
        \item Local cost functions (support on $O(1)$ qubits)
        \item Problem-inspired ansätze (exploit structure)
        \item Layer-wise training (incremental depth)
        \item Limiting entanglement in early training
      \end{itemize}
    \column{0.44\textwidth}
      \begin{block}{Entanglement regimes}
        \begin{tabular}{ll}
          \toprule
          Entanglement & Training \\
          \midrule
          Low & Classically simulable \\
          Moderate & Optimal learning \\
          High & Barren plateau \\
          \bottomrule
        \end{tabular}
      \end{block}
      \vspace{0.3cm}
      \begin{alertblock}{Quantum Neural Tangent Kernel}
        In the linearised training regime:
        \[
          K(x,x')=\sum_i
            \frac{\partial f(x)}{\partial\theta_i}
            \frac{\partial f(x')}{\partial\theta_i}
        \]
        connects QNNs to kernel methods.
      \end{alertblock}
  \end{columns}
\end{frame}

%% ── Variational classifier circuit ──────────────────────────────
\begin{frame}{Variational Quantum Classifier: Full Circuit}
  Three-stage circuit for binary classification on $n=3$ qubits:

  \vspace{0.3cm}
  \begin{center}
  \begin{quantikz}[row sep=0.5cm, column sep=0.4cm]
    \lstick{$\ket{0}$}
      & \gate{R_y(x_1)}
      & \gate{R_z(\alpha_1)} & \ctrl{1} & \qw
      & \gate{R_z(\beta_1)} & \ctrl{1} & \qw
      & \meter{$Z$} \\
    \lstick{$\ket{0}$}
      & \gate{R_y(x_2)}
      & \gate{R_z(\alpha_2)} & \targ{} & \ctrl{1}
      & \gate{R_z(\beta_2)} & \targ{} & \ctrl{1}
      & \qw \\
    \lstick{$\ket{0}$}
      & \gate{R_y(x_3)}
      & \gate{R_z(\alpha_3)} & \qw & \targ{}
      & \gate{R_z(\beta_3)} & \qw & \targ{}
      & \qw
  \end{quantikz}
  \end{center}

  \vspace{0.2cm}
  \begin{itemize}
    \item \textbf{Encoding} (col.~2): $R_y(x_i)$ angle-embeds each feature
    \item \textbf{Var.\ layer 1} (cols.~3--5): $R_z$ rotations + CNOT chain $0\to1\to2$
    \item \textbf{Var.\ layer 2} (cols.~6--8): second block of $R_z$ + CNOT chain
    \item \textbf{Output}: $f = \langle Z_0\rangle$; class~1 if $f>0$, class~0 if $f<0$
  \end{itemize}
\end{frame}

%% ── Entanglement and expressivity ───────────────────────────────
\begin{frame}{Entanglement Entropy and Expressivity}
  \begin{columns}[T]
    \column{0.52\textwidth}
      Entanglement entropy of subsystem $A$:
      \[
        S(\rho_A) = -\mathrm{Tr}(\rho_A\log\rho_A)
      \]
      Controls the complexity of the variational state:
      \begin{itemize}
        \item $S=0$: product state, no correlations, classically
              efficient
        \item $S = O(\text{area law})$: MPS-representable
        \item $S = O(n)$: volume-law, generic circuit output
      \end{itemize}

      \medskip
      The QNN ansatz lives on a \emph{low-dimensional submanifold}
      of the full exponentially-large Hilbert space.
      The manifold dimension is $\dim\theta$.
    \column{0.44\textwidth}
      \begin{alertblock}{Expressivity bound}
        A circuit with $L$ layers and $n$ qubits can generate states
        with at most
        \[
          S(\rho_A) \lesssim L\log 2
        \]
        entanglement entropy between any bipartition — a property
        used to prove simulability of shallow circuits.
      \end{alertblock}
      \vspace{0.2cm}
      \begin{block}{Optimal regime}
        Moderate $S$ (moderate entanglement): rich correlations
        without concentration of measure.
      \end{block}
  \end{columns}
\end{frame}

%% ── PyTorch simulation summary ───────────────────────────────────
\begin{frame}[fragile]{Simulating a QNN: NumPy + PyTorch}
  We simulate the circuit as $2^n\times 2^n$ unitary matrix
  operations in NumPy; gradients use the parameter-shift rule;
  optimisation uses PyTorch Adam.

  \medskip
  \begin{lstlisting}
def var_layer(p):
    """Ry(p0) x Ry(p1) -> CNOT -> Rz(p2) x Rz(p3)."""
    return kron(rz(p[2]),rz(p[3])) @ CNOT @ kron(ry(p[0]),ry(p[1]))

def circuit(params, x):
    """Full 2-qubit VQC: Rx feature map + 2 variational layers."""
    state = |00>  # complex vector, length 4
    U = kron(Rx(x[0]), Rx(x[1]))          # feature map
    U = var_layer(params[:4]) @ U          # layer 1
    U = var_layer(params[4:8]) @ U         # layer 2
    psi = U @ state
    return <psi| Z0 |psi>                  # expectation value

def grad_circuit(params, x, shift=pi/2):
    """Exact gradient via parameter-shift rule."""
    g = zeros(8)
    for i in range(8):
        g[i] = 0.5*(circuit(params + shift*e_i, x)
                  - circuit(params - shift*e_i, x))
    return g
  \end{lstlisting}
\end{frame}

%% ── 3-qubit credit classifier circuit ───────────────────────────
\begin{frame}{Credit Classifier: 3-Qubit Strongly-Entangling Circuit}
  \begin{center}
  \begin{quantikz}[row sep=0.5cm, column sep=0.4cm]
    \lstick{$\ket{0}$}
      & \gate{R_y(x_1)}
      & \gate{R_z R_y R_z(\alpha_1,\beta_1,\gamma_1)}
      & \ctrl{1} & \qw
      & \meter{$Z$} \\
    \lstick{$\ket{0}$}
      & \gate{R_y(x_2)}
      & \gate{R_z R_y R_z(\alpha_2,\beta_2,\gamma_2)}
      & \targ{} & \ctrl{1}
      & \qw \\
    \lstick{$\ket{0}$}
      & \gate{R_y(x_3)}
      & \gate{R_z R_y R_z(\alpha_3,\beta_3,\gamma_3)}
      & \ctrl{-2} & \targ{}
      & \qw
  \end{quantikz}
  \end{center}

  \vspace{0.3cm}
  \begin{columns}[T]
    \column{0.48\textwidth}
      \textbf{Encoding:} $R_Y(x_i)$ angle-embeds income, debt ratio,
      age (scaled to $[0,\pi]$). Each feature becomes a rotation angle
      on one qubit.
    \column{0.48\textwidth}
      \textbf{Ansatz:} general $SU(2)$ Euler rotation per qubit
      ($R_z R_y R_z$ = 3 parameters per qubit),
      then \emph{ring} of CNOTs ($0\!\to\!1\!\to\!2\!\to\!0$).
      Total: 9 trainable parameters.
  \end{columns}
\end{frame}

%% ============================================================
\section{Variational Quantum Eigensolver}
%% ============================================================

\begin{frame}{Variational Quantum Eigensolver (VQE)}
  \begin{columns}[T]
    \column{0.52\textwidth}
      \textbf{Variational principle:}
      \[
        E_0 = \min_\psi\langle\psi|H|\psi\rangle,\quad
        E(\theta)\ge E_0
      \]
      \textbf{Ansatz:} $\ket{\psi(\theta)} = U(\theta)\ket{0}$

      \textbf{Measurement:} decompose $H = \sum_i c_i P_i$:
      \[
        E(\theta) = \sum_i c_i \langle P_i\rangle
      \]
      Each $\langle P_i\rangle$ measured on the quantum device.

      \textbf{Gradient} via parameter-shift rule:
      \[
        \frac{\partial E}{\partial\theta}
        = \tfrac{1}{2}\left[E(\theta+\tfrac{\pi}{2})
                           -E(\theta-\tfrac{\pi}{2})\right]
      \]
    \column{0.44\textwidth}
      \textbf{UCC ansatz:}
      \[
        U = e^{T-T^\dagger},\quad
        T = \sum_{ai}t_{ai}\,a_a^\dagger a_i + \cdots
      \]
      Trotterised as $e^{A+B}\approx e^Ae^B$.

      BCH connection to coupled-cluster:
      \[
        e^{-T}He^T = H + [H,T] + \tfrac{1}{2}[[H,T],T]+\cdots
      \]

      \begin{center}
      \begin{quantikz}[row sep=0.4cm, column sep=0.4cm]
        \lstick{$\ket{0}$}
          & \gate[4,nwires=3]{U(\theta)} & \meter{} \\
        \lstick{$\ket{0}$} & & \meter{} \\
        \lstick{$\ket{0}$} & & \qw \\
        \lstick{$\ket{0}$} & \qw & \qw
      \end{quantikz}
      \end{center}
  \end{columns}
\end{frame}

%% ============================================================
\section{Quantum Approximate Optimization Algorithm}
%% ============================================================

\begin{frame}{QAOA: Problem Statement and Ansatz}
  \textbf{Goal:} maximise a classical cost $C(z)$ over bit strings
  $z\in\{0,1\}^n$.

  Encode as a diagonal Hamiltonian: $H_C\ket{z} = C(z)\ket{z}$.

  \medskip
  \textbf{QAOA ansatz} (depth $p$):
  \[
    \ket{\gamma,\beta}
    = \prod_{l=1}^p e^{-i\beta_l H_M}\,e^{-i\gamma_l H_C}\,\ket{+}^{\otimes n}
  \]
  \begin{columns}[T]
    \column{0.48\textwidth}
      \begin{center}
      \begin{quantikz}[row sep=0.42cm, column sep=0.3cm]
        \lstick{$\ket{+}$}
          & \gate[3,nwires=2]{e^{-i\gamma H_C}}
          & \gate{R_x(2\beta)} & \qw \\
        \lstick{$\ket{+}$}
          &  & \gate{R_x(2\beta)} & \qw \\
        \lstick{$\ket{+}$}
          & \qw & \gate{R_x(2\beta)} & \qw
      \end{quantikz}
      \end{center}
      {\small One QAOA layer ($p=1$): cost unitary then mixer.}
    \column{0.48\textwidth}
      \begin{itemize}
        \item $H_M = \sum_i X_i$: mixer Hamiltonian
        \item $\ket{+}=H\ket{0}$: uniform superposition
        \item Maximise $C(\gamma,\beta)=\langle H_C\rangle$ over
              $2p$ parameters
      \end{itemize}
      \begin{alertblock}{Adiabatic connection}
        QAOA = Trotterised adiabatic evolution.
        $p\to\infty$ recovers exact adiabatic QC.
      \end{alertblock}
  \end{columns}
\end{frame}

\begin{frame}{QAOA: MaxCut and Trotterisation}
  \textbf{MaxCut} on graph $G=(V,E)$: cost Hamiltonian
  \[
    H_C = \sum_{\langle ij\rangle\in E}\frac{1-Z_iZ_j}{2}
  \]

  The cost unitary $e^{-i\gamma H_C}$ decomposes into two-qubit
  $R_{zz}$ gates:
  \begin{center}
  \begin{quantikz}[row sep=0.5cm, column sep=0.45cm]
    \lstick{$\ket{+}$} & \gate{H} &
      \ctrl{1} & \gate{R_z(2\gamma)} & \ctrl{1} & \gate{R_x(2\beta)} &
      \meter{} \\
    \lstick{$\ket{+}$} & \gate{H} &
      \targ{} & \qw & \targ{} & \gate{R_x(2\beta)} &
      \qw
  \end{quantikz}
  \end{center}
  {\small $p=1$ QAOA for a single edge $(i,j)$: CNOT–$R_z$–CNOT implements
   $e^{-i\gamma Z_iZ_j/2}$ on qubits $i,j$; $R_x(2\beta)$ is the mixer.}

  \medskip
  \begin{alertblock}{Trotter expansion}
    \[
      e^{-i(H_M+H_C)t}\approx e^{-iH_M\Delta t}e^{-iH_C\Delta t}
    \]
    QAOA at depth $p$ = $p$-step Trotterised adiabatic evolution with
    variable step sizes $(\gamma_l,\beta_l)$.
  \end{alertblock}
\end{frame}

\begin{frame}{QAOA: Gradient, Landscape, and Relations to QNN/VQE}
  \begin{columns}[T]
    \column{0.52\textwidth}
      \textbf{Gradient} w.r.t.\ $\gamma_l$:
      \[
        \partial_{\gamma_l}C
        = i\langle[H_C,H_{\text{eff}}]\rangle
      \]
      Computed via \emph{finite differences} (not parameter-shift,
      because $H_C$ has a multi-valued spectrum $\{0,\ldots,|E|\}$).

      \medskip
      \textbf{Landscape:} $C(\gamma,\beta)$ non-convex but has
      structured correlations — warm-starting and parameter
      concentration exploit this.
    \column{0.44\textwidth}
      \begin{tabular}{lcc}
        \toprule
        Feature & QAOA & VQE \\
        \midrule
        Ansatz & fixed & general \\
        Motivation & adiabatic & variational \\
        Target & opt.\ problem & ground state \\
        Params & $2p$ & many \\
        \bottomrule
      \end{tabular}

      \vspace{0.3cm}
      \begin{block}{QAOA $\subset$ QNN}
        QAOA is a QNN with \emph{specific} Hamiltonians $H_C$, $H_M$
        instead of a general hardware-efficient ansatz.
      \end{block}
  \end{columns}
\end{frame}

%% ============================================================
\section{Quantum Support Vector Machines}
%% ============================================================

\begin{frame}{QSVM: Quantum Kernel Methods}
  \begin{columns}[T]
    \column{0.52\textwidth}
      \textbf{Quantum feature map:}
      \[
        \ket{\phi(x)} = U(x)\ket{0},\quad
        K(x,x') = |\langle\phi(x)|\phi(x')\rangle|^2
      \]
      The kernel is a \textbf{fidelity} — a quantum overlap, directly
      analogous to many-body correlation functions.

      \medskip
      \textbf{Swap test circuit:}
      \begin{center}
      \begin{quantikz}[row sep=0.4cm, column sep=0.4cm]
        \lstick{$\ket{0}$} & \gate{H} & \ctrl{1} & \gate{H} & \meter{} \\
        \lstick{$\ket{\phi(x)}$} & \qw & \swap{1} & \qw & \qw \\
        \lstick{$\ket{\phi(x')}$} & \qw & \targX{} & \qw & \qw
      \end{quantikz}
      \end{center}
      $P(0) = \tfrac{1}{2}(1+K(x,x'))$
    \column{0.44\textwidth}
      \begin{block}{SVM dual optimisation}
        \[
          \max_\alpha\sum_i\alpha_i
          - \tfrac{1}{2}\sum_{ij}\alpha_i\alpha_j y_iy_j K_{ij}
        \]
        Solved classically once $K_{ij}$ is estimated on the device.
      \end{block}
      \vspace{0.2cm}
      \begin{alertblock}{Concentration risk}
        High entanglement in $U(x)$:
        \[
          K(x,x')\approx\frac{1}{2^n}\;\forall\,x,x'
        \]
        Kernel becomes uninformative — avoid deep random feature maps.
      \end{alertblock}
      \vspace{0.2cm}
      \textbf{Connection to HHL:} kernel system
      $(K+\lambda I)\alpha=y$ solvable by HHL.
  \end{columns}
\end{frame}

%% ============================================================
\section{HHL Algorithm}
%% ============================================================

\begin{frame}{HHL: Quantum Linear Systems Solver}
  \begin{columns}[T]
    \column{0.52\textwidth}
      \textbf{Problem:} given Hermitian $A$ and $\ket{b}$, prepare
      $\ket{x}\propto A^{-1}\ket{b}$.

      \textbf{Spectral decomposition:}
      \[
        A^{-1}\ket{b}
        = \sum_j\frac{\beta_j}{\lambda_j}\ket{u_j}
      \]
      HHL implements $\lambda_j\to\lambda_j^{-1}$ via QPE.

      \textbf{Steps:}
      \begin{enumerate}
        \item QPE: encode eigenvalues in register
        \item Controlled rotation: apply $\lambda_j^{-1}$
        \item Uncompute QPE; post-select ancilla $=\ket{1}$
      \end{enumerate}

      \textbf{Complexity:}
      $\mathcal{O}(\kappa^2\log N)$ vs classical $O(N\sqrt\kappa)$.
    \column{0.44\textwidth}
      \begin{center}
      \begin{quantikz}[row sep=0.35cm, column sep=0.3cm]
        \lstick{$\ket{b}$}
          & \gate[1]{QPE} & \qw
          & \gate[1]{QPE^\dagger} & \qw \\
        \lstick{$\ket{0}^{\otimes m}$}
          & \qw & \gate{\lambda^{-1}} & \qw & \qw \\
        \lstick{$\ket{0}$\,(anc.)}
          & \qw & \ctrl{-1} & \qw & \meter{}
      \end{quantikz}
      \end{center}
      \vspace{0.2cm}
      \begin{alertblock}{Green's function connection}
        \[
          A^{-1}\sim(\omega I-H)^{-1} = G(\omega)
        \]
        HHL $\equiv$ quantum computation of the resolvent operator /
        Green's function of the many-body Hamiltonian.
      \end{alertblock}
  \end{columns}
\end{frame}

%% ============================================================
\section{Grand Unification: QML as Operator Science}
%% ============================================================

\begin{frame}{Three Fundamental Structures of QML}
  \[
    \boxed{
      \text{Quantum Machine Learning}
      = \text{Geometry} + \text{Dynamics} + \text{Operators}
    }
  \]
  \vspace{0.2cm}
  \begin{center}
  \renewcommand{\arraystretch}{1.3}
  \begin{tabular}{p{2.4cm}p{4.2cm}p{4.0cm}}
    \toprule
    \textbf{Structure} & \textbf{Mathematical object} & \textbf{QML method} \\
    \midrule
    Geometry  & QFI metric $F_{ij}$, Hilbert space $\mathcal{H}$
              & QSVM (kernel = fidelity) \\
    Dynamics  & Variational evolution, TDVP
              & QNN, VQE, QAOA \\
    Operators & Inverse maps, Green's functions
              & HHL, QPE \\
    \bottomrule
  \end{tabular}
  \end{center}

  \vspace{0.3cm}
  \textbf{Unifying correspondences:}
  \begin{align*}
    \text{HHL} &\;\longleftrightarrow\; \text{exact inverse operator}\\
    \text{QSVM} &\;\longleftrightarrow\; \text{kernel geometry }
      K(x,x')=|\langle\phi(x)|\phi(x')\rangle|^2\\
    \text{QNN/VQE/QAOA} &\;\longleftrightarrow\; \text{variational dynamics }
      \textstyle\sum_j F_{ij}\dot\theta_j = C_i
  \end{align*}
\end{frame}

\begin{frame}{Spectral Perspective and Learning as Projection}
  \begin{columns}[T]
    \column{0.52\textwidth}
      All methods rely on the spectral structure:
      \[
        A = \sum_j\lambda_j\ket{u_j}\bra{u_j}
      \]
      \begin{itemize}
        \item \textbf{HHL:} transforms $\lambda_j\to\lambda_j^{-1}$
        \item \textbf{QNN/VQE:} learn the eigenstructure via
              gradient descent on the variational manifold
        \item \textbf{QSVM:} measures pairwise overlaps between
              eigenstates
      \end{itemize}

      \medskip
      \textbf{Learning as projection:}
      \[
        \text{Learning}\equiv\text{optimal projection onto}
      \]
      \[
        \text{a submanifold of Hilbert space}
      \]
    \column{0.44\textwidth}
      \begin{alertblock}{Unified master equation}
        \[
          \boxed{
            \text{Learning}\equiv
            \text{approximating operators}
          }
        \]
      \end{alertblock}
      \vspace{0.3cm}
      \begin{tabular}{ll}
        \toprule
        Approach & Operator \\
        \midrule
        Exact & HHL ($A^{-1}$) \\
        Variational & QNN/VQE/QAOA \\
        Geometric & QSVM \\
        \bottomrule
      \end{tabular}
      \vspace{0.3cm}

      \textbf{Information geometry:}
      $ds^2=\sum_{ij}F_{ij}d\theta_id\theta_j$;
      training follows geodesics of the QFI metric.
  \end{columns}
\end{frame}

%% ============================================================
\section{Outlook and References}
%% ============================================================

\begin{frame}{Outlook: Quantum Advantage and Future Directions}
  \begin{columns}[T]
    \column{0.50\textwidth}
      \begin{block}{Where can genuine advantage emerge?}
        \begin{itemize}
          \item \textbf{Fault-tolerant QC:} HHL, QPE, Shor —
                provable exponential speedup
          \item \textbf{NISQ VQA:} VQE, QAOA, QSVM —
                heuristic, unclear general advantage
          \item \textbf{Quantum kernels:} only when the feature map
                is classically hard to simulate
        \end{itemize}
      \end{block}
      \vspace{0.3cm}
      \begin{alertblock}{Fundamental trade-off}
        \[
          \underbrace{\text{High expressivity}}_{\text{deep circuits}}
          \;\longleftrightarrow\;
          \underbrace{\text{Hard optimisation}}_{\text{barren plateaus}}
        \]
      \end{alertblock}
    \column{0.46\textwidth}
      \textbf{Future directions:}
      \begin{itemize}
        \item Fault-tolerant QC and error correction
        \item Adaptive ansatz design (ADAPT-VQE)
        \item Quantum-classical co-design (hardware-aware)
        \item Connection to many-body physics:
              VQE $\leftrightarrow$ coupled cluster,
              QAOA $\leftrightarrow$ adiabatic,
              HHL $\leftrightarrow$ Green's functions
        \item Quantum transformers and quantum RL
        \item Field-theoretic view: circuits as discretised
              quantum field theories
      \end{itemize}
  \end{columns}
\end{frame}

\begin{frame}{References}
  \begin{columns}[T]
    \column{0.50\textwidth}
      \textbf{QNN and VQC}
      \begin{enumerate}
        \item A.~Abbas et al., \emph{Power of QNNs}, Nature CS \textbf{1}, 403 (2021).
        \item E.~Anschuetz \& B.~Kiani, Nat.\ Commun.\ \textbf{13}, 7760 (2022).
        \item M.~Zhao et al., arXiv:2504.16131 (2025).
      \end{enumerate}
      \textbf{QAOA}
      \begin{enumerate}\setcounter{enumi}{3}
        \item E.~Farhi et al., arXiv:1411.4028 (2014).
        \item K.~Blekos et al., Phys.\ Rep.\ \textbf{1029}, 1 (2024).
        \item L.~Zhou et al., Phys.\ Rev.\ X \textbf{10}, 021067 (2020).
      \end{enumerate}
    \column{0.46\textwidth}
      \textbf{VQE}
      \begin{enumerate}\setcounter{enumi}{6}
        \item A.~Peruzzo et al., Nat.\ Commun.\ \textbf{5}, 4213 (2014).
        \item J.~Tilly et al., Phys.\ Rep.\ \textbf{986}, 1 (2022).
      \end{enumerate}
      \textbf{QSVM}
      \begin{enumerate}\setcounter{enumi}{8}
        \item V.~Havlíček et al., Nature \textbf{567}, 209 (2019).
        \item M.~Schuld \& N.~Killoran, PRL \textbf{122}, 040504 (2019).
      \end{enumerate}
      \textbf{HHL and QFI}
      \begin{enumerate}\setcounter{enumi}{10}
        \item A.~W.~Harrow et al., PRL \textbf{103}, 150502 (2009).
        \item J.~Stokes et al., Quantum \textbf{4}, 269 (2020).
        \item J.~R.~McClean et al., Nat.\ Commun.\ \textbf{9}, 4812 (2018).
      \end{enumerate}
  \end{columns}
\end{frame}

\end{document}
